% !TEX root = ../../main.tex
\documentclass[../../main.tex]{subfiles}

\begin{document}
\title{
  Galois Groups of Trinomials with Small Coefficients
}
\author{Hilel Garmi}
\maketitle

Hello everyone, I am Hilel Garmi, and I will talk about Galois groups of trinomials with small coefficients.
Our main example is going to be the family of polynomials of the form [frame]
\[
  f = x^n - 2x + 1
\]
I will first try to convince why it is interesting.
[frame]
The first reason is that in 1892 Hilbert proved his Irreducibility Theorem, which implies that [slide] in some sense, "almost all" polynomials of degree $n$ have Galois group $S_n$ over $\Q$.
However, contrary to what one might expect upon hearing about this theorem for the first time, it turns out to be very hard to give explicit examples of a polynomial of degree $n$ with Galois group $S_n$ for every $n$.
Almost 40 years later, Schur found the first example [slide] of such a family:
\[
  f := \sum_{k=0}^n \binom{n}{k} \frac{x^k}{k!}
\]
My problem with this example, is that both the numerator and the denominator of the coefficients grow insanely fast with $n$, and I would like to be able to find examples whose coefficients are (1) integers, and (2) bounded.
And while until recently it might have seemed like an arbitrary challenge to choose, in 2020 Lior Bary-Soroker from Tel Aviv and Gady Kozma from the Weizmann proved that [slide] if you sample the coefficients of a monic polynomial of degree $n$ uniformly and independently from $\left\{ 1,\dots,210 \right\}$ then the probability that its Galois group is either $A_n$ or $S_n$ tends to $1$ as $n$ tends to infinity, and they conjecture that the same is true for $S_n$.
So just as Hilbert's theorem implied that we should be able to find many examples as the one that Schur eventually provided, Bary-Soroker and Kozma's result makes it a bit embarrassing for number theorists that it seems like the only known example with bounded integer coefficients is [slide] 
\[
  f = x^n - x - 1
\]
which was proven to be irreducible by Selmer in 1956, and to have Galois group $S_n$ by Nart and Vila in 1979.
So our first motivation is just ease this embarrassment a bit.

% [frame]
% The second motivation is the embaressment arising from the fact that we know almost nothing about how to calculate Galois groups of polynomials over $\Q$. I mean [slide] Kummer theory solves completely the case $x^n + b$, but once we try to consider trinomials of the form [slide] $x^n + ax^k + b$, apart from calculating it using a computer if $n$ is small enough we are pretty much lost. 
% So if we will be able to solve the case $x^n - 2x + 1$, and generalize it a bit further, it can maybe be a first step towards developing a more satisfactory theory for Galois groups of trinomials.

Before, I will get to talk about $x^n -2x + 1$, I will sketch the proof of the case $x^n - x - 1$, which will be instructive for us.
To do that, I will first need to tell you about inertia groups.
[frame]

First a piece of notation, for any field $\F$, let $\overline{\F}$ denote its algebraic closure, and $G_\F := \Gal\left( \overline{\F}/\F \right)$ denote its absolute Galois group.
Secondly, just like $\Z_p$ [slide] is a local ring with maximal ideal $p\Z_p$ and residue field $\F_p$, the same image appears in each algebraic extension $L/\Q_p$.
The integral closure $\mathcal{O}_L$ of $\Z_p$ inside of $L$ is again a local ring, and its residue field is some algebraic extension of $\F_p$.
The same is true even when $L$ is the algebraic closure $\overline{\Q_p}$ of $\Q_p$, and in this case the residue field is the algebraic closure $\overline{\F_p}$ of $\F_p$.
Now it turns out that [slide] any automorphism in $\overline{\Q_p}$ over $\Q_p$ induces an automorphism of $\overline{\F_p}$.
[slide]
Thus we have a map $G_{\Q_p} \to G_{\F_p}$.
The absolute inertia group $I_{\Q_p}$ [slide] is the kernel of this map, i.e. the group of automorphisms of $\overline{\Q_p}$ over $\Q_p$ that induce the identity on $\overline{\F_p}$.
So in this picture, each of the automorphisms in $I_{\Q_p}$ move each element of $\overline{\Q_p}$ into some element within the same circle.

Now we may view the roots of $f$ as elements of $\overline{\Q_p}$, and ask what we can say about the action of $I_{\Q_p}$ on them.
Well, by the definition we just gave, for two roots to be in the same orbit of $I_{\Q_p}$, they must live in the same circle.
Hence, this circle represnts an element of $\overline{\F_p}$, that is a double root of $f$.
If three roots are in the same orbit, then this circle represents a triple root of $f$ and so forth. 
If we derive $f$ twice we get $f'' = n\left( n-1 \right) x^{n-2}$.
And well, at least if $p$ does not divide $n\left( n-1 \right)$, then the only root of $f''$ is $0$, which is not a root of $f$, and hence $f$ does not have a triple root.
This is just to give you a taste of the kind of arguments that go into the proof.
It turns out that we can bypass the problem of $p$ dividing $n\left( n-1 \right)$, and show that for any prime $p$, $f$ has no triple roots, and in fact has at most one double root.
Therefore, there is at most one pair of roots of $f$ that live in the same circle, and hence $I_{\Q_p}$ acts on the roots of $f$ either trivially, or as a single transposition, transposing these two roots that live in the same circle.

Now all of this is related to the Galois group of our polynomial over $\Q$?
[slide]
Well inside $\overline{\Q_p}$, there is a copy of $\overline{\Q}$, and hence we can reduce any automorphism in $I_{\Q_p}$ to an automorphism of $\overline{\Q}$.
Thus, we have an image of $I_{\Q_p}$ inside $G_\Q$, and since it was obtained from $I_{\Q_p}$, just by reducing to $\overline{\Q}$, it acts on the roots of $f$ in the same way as $I_{\Q_p}$ does, i.e. either trivially or as a single transposition.
Note by the way, that this upper arrow is not cannonic at all, since we can precompose it with any of the many automorphisms of $\overline{\Q}$.
Each such choice of embdding is called a place above $p$ for reasons I will not discuss, so each place above $p$ induces another embedding of $I_{\Q_p}$ into $G_\Q$ which is called the inertia at this place.
We can now utilize a theorem which is called Hermite's Theorem, which says, certainly not in its most well known form, that the Galois group of a polynomial over $\Q$ is generated by its inertia groups.
In our case, we just saw that each inertia group acts on the roots of $f$ either trivially or as a single transposition, and hence the Galois group of $f$ over $\Q$ is generated by transpositions.

Finally, we can speak about $x^n -2x + 1$.  
[frame]
The first reason that I can not show you that the Galois group of $f = x^n - 2x + 1$ is $S_n$, is that it is not $S_n$, and for every obvious reason: 
Its Galois group is not transitive, because $f$ is reducible:
Indeed $1$ is a root of $f$, and hence $f$ factors as [slide]
\[
  x^n - 2x + 1 = \left( x-1 \right) g 
\]
where 
\[
  g = x^{n-1} + x^{n-2} + \cdots + x - 1
\]
So the biggest the Galois group can be is $S_{n-1}$.
But never mind, both motivations I stated at the beginning still apply, $g$ will give us an example of a polynomial with bounded integer coefficients and Galois group $S_{n-1}$, and we will have one more trinomial whose Galois group we can calculate.
$g$ attracted its own share of attention, because it is kind of another disguise of the [slide] generalized Fibonacci polynomial 
\[
  x^{n-1} - x^{n-2} - \cdots - x - 1
\]
which was proven to be irreducible by Miles in 1960, and more explicitly Wolfram in 1998.
The irreducibility of $g$ can then be easily deduced.
We are also not the first to try to calculate its Galois group and in 2003 Martin succeeded to show that [slide] it is $S_{n-1}$ for any odd value of $n$.
What he did is essentially running the same argument that I just sketched for $x^n - x - 1$.
The fact that $x^n - 2x + 1$ is reducible, does not hurt the argument, since the roots of $g$ are all roots of $f$ and hence if $f$ has at most on double root above each prime, then so does $g$, and also all the details that I hid from you worked just fine.
The reason it fails for even values of $n$, is that in this case
\[
  x^n - 2x + 1 \equiv x^n - 1 \equiv \left( x^{n/2} - 1\right)^2 \mod 2
\]
and hence every root of $f$ in $\overline{\F_2}$ is a double root, and hence, all the roots of $g$ in $\overline{\F_2}$, well except $1$, are double roots, and thus nothing prevents a membed of the inertia group over $\Q_2$ from tranposing for example the two roots in this circle at the same time.
So, we do not know that the Galois group of $g$ over $\Q$ is generated by transpositions anymore. 

But do not you worry, we found a remedy.
Over all the primes $p \neq 2$, the inertia group at every place still acts on the roots of $g$ either trivially or as a single transposition.
Let then $N$ denote the subgroup of $G_\Q$ generated by the inertia groups at all the places over all the primes $p \neq 2$. 
group of $g$ generated by the inertia groups at all of these places.
If we show that $N$ is transitive, then by the same fact we used earlier, $N$ is $S_{n-1}$, and hence the Galois group of $g$ over $\Q$ is $S_{n-1}$ as well.

Now we are going to do a slight change of perspective.
Let $L = \Q\left[ x \right] / (g)$ be the extension of $\Q$ defined by $g$.
By the universal properties of $\Q\left[ x \right]$ and of $g$, roots of $g$ in $\overline{\Q}$ are like embeddings of $L$ into $\overline{\Q}$, and the action of the Galois group of $g$ over $\Q$ on its roots is like the action of $G_\Q$ on these embeddings.
The reason I am doing this change of perspectives, is to utilize the fact that the functor, $\Hom_\Q\left( -, \overline{\Q} \right)$, is a contravariant equivalence between the category of finite extensions of $\Q$ and finite transitive $G_\Q$-sets.  

Now assume this is the set of embeddings of $L$ into $\overline{\Q}$, we can divide it into orbits of $N$.
As its name suggests, $N$ turns out to be a normal subgroup of the Galois group of $g$, and hence the action of $G$ on $\Hom\left( L, \overline{\Q} \right)$ respect the divison into orbits of $N$, and hence induces an action on the set of orbits of $N$.
Therefore, we got a $G_\Q$-equivariant map from $\Hom\left( L, \overline{\Q} \right)$ to another set $\Hom\left( L, \overline{\Q} \right) / N$.
Therefore, we have a corresponding morphism in the opposite direction in the category of finite extensions of $\Q$, i.e. some corresponding intermediate extension $K/ \Q$.
If we assume towards contradiction that $N$ is not transitive, then the quotint by it consists of more than one element, and hence $K$ has more than one embedding into $\overline{\Q}$, and hence $K$ is not $\Q$.
Therefore, it suffices for us to show that such an intermediate extension $K$ does not exist.

However, note that in this picture, each element of each inertia group at a place over a prime $p \neq 2$, act on the members of this set, either trivially or by transposition.
Therefore, the two elments it transposes must be in the same orbit of $N$, otherwise it would have to move more than two elements around.
Therefore, the action that it induces on the set of orbits of $N$ is trivial, and hence its action on $K$ is trivial as well.
Therefore, $K$ is fixed by all inertia groups, except those that are over $2$.
Now I do not have time to talk about it, but it turns out that this condition is pretty restrictive, and by proving some bounds on the discriminant of $L$, for those of you who know what this means, we were able also bound the root discriminant of $K$. 
Combining this fact with known listings of number fields with small root discriminant, we were able to show that such a $K$ does not exist, from which it follows that $N$ is transitive, from which it follows that the Galois group of $g$ over $\Q$ is $S_{n-1}$.



No one said we need to show that the entire Galois group of $g$ is generated by a transpositions, we only need to show that some transitive subgroup of it is generated by transpositions, and at every place over any prime $p \neq 2$ the inertia group does act on $g$ is either trivially or as a single transposition.  
Let then $N$ denote the subgroup of the Galois group of $g$ over $\Q$ that is generated by the inertia groups of all the places over all the primes $p \neq 2$.




% and a lot of time it turns out, it is easier to understand the action of $G_{\Q_p}$ on the roots of a polynomial than the action of $G_\Q$ on them.
% If we have done so for different primes $p$, we can then combine the information we have about the action of $G_{\Q_p}$ on the roots of a polynomial to get information about the action of $G_\Q$ on them.
% [slide]
% What Schur did in his proof, was to show that there is a prime $\frac{n}{2} < p < n - 2$ such that the Galois group of $f$ over $\Q_p$ contains a $p$-cycle.
% Hence the action of $G_\Q$ on the roots of his $f$ is transitive, because one can show that  


% The first part of the argument is called the local part, and the second part is called the global part.
% So 
% In particular, any automorphism in $I_{\Q_p}$ does so, and it might be easier   


% When we ask what is the Galois group of a polynomial $f$ as a subgroup of $S_n$, we essentially ask what are all the permutations of its roots that are induced by automorphisms of $\overline{\Q}$.



% Firs of all, [slide] I can view the $p$-adic integers $\Z_p$ as an infinite $p$-regular tree, in which each vertex of depth $d$ consists of all the $p$-adic integers that are equivalent modulo $p^d$.

% Now it turns out, that the same image appears in any algebraic extension $L/\Q_p$. 
% The integral closure [slide] $\mathcal{O}_L$ of $\Z_p$ inside of $L$ is again a local ring, and if we denote its maximal ideal by $\p_L$, we can view $\mathcal{O}_L$ as an infinite tree, [slide] in which each vertex of depth $d$ consists of all the elements of $\mathcal{O}_L$ that are equivalent modulo $\p_L^d$.
% Where earlier the children of each vertex were indexed by $\F_p$, now they are indexed by the residue field $\mathcal{O}_L/\p_L$ of $L$, which is some field extension of $\F_p$.
% As a warning I will say that sometimes there are also more intermediate levels between each two levels corresponding to two consecutive powers of $p$ in the first tree, but we are not going to talk about it.

% The same is true even when $L$ is the algebraic closure $\overline{\Q_p}$ of $\Q_p$, and in this case the residue field is the algebraic closure $\overline{\F_p}$ of $\F_p$.

% Now any automorphism of $\overline{\Q_p}$ over $\Q_p$ preserves this tree structure, and we can define the absolute inertia group $I_{\Q_p}$ to be the subgroup of $\Gal\left( \overline{\Q_p}/\Q_p \right)$ that fixes the first level of the tree, which is just $\overline{\F_p}$.
% In an algebraic language, $I_{\Q_p}$ is the subgroup of 

% $\Z_p$ is a local ring with maximal ideal $p\Z_p$ and residue field $\F_p$, and I can use this fact to define the weird $p$-adic metric, which assign to each pair of $p$-adic numbers a distance out of 

% First of all note that if we choose an algebraic closure $\overline{\Q_p}$ of $\Q_p$, it contains an algebraic closure of $\Q$, so we can choose an embedding that makes [slide] this diagram commute.


\end{document}